\chapter{EXUMAÇÃO, DIÁRIO E CONTAMINAÇÃO}

\epigraph{O senhor vai me chamar de louco. O senhor já está me chamando. Mas eu vi. O olho brilhava. O caderno estava na mão dele como se ele tivesse esperando a gente o tempo inteiro. Ele sabia. Ele sabia que a gente ia abrir aquela cova. E ele sabia que quem lesse o caderno ia continuar o que ele começou.}{--- Genésio Barbosa, coveiro do Hospital Colônia (depoimento à Dra. Lúcia Menezes, 17/06/1979)}

\vspace{1em}

\section*{4. ACHADOS ESPECIAIS --- EXUMAÇÃO}

\noindent\textit{14 de junho de 1979. 48 horas após o óbito. Autorização verbal da direção do Hospital Colônia para liberação do leito. Testemunhas: Dr. Álvaro Torres (legista), Genésio Barbosa (coveiro), funcionário administrativo não identificado.}
\vspace{1em}

\subsection*{4.1 Depoimento de Genésio Barbosa}

\noindent\textit{Transcrição literal do depoimento gravado. Genésio trabalha no Colônia há 22 anos. Durante o depoimento, segurava o chapéu com as duas mãos e não olhava para o gravador.}
\vspace{1em}

\lettrine{E}{}u já enterrei muito homem. Muita mulher. Muita criança. O Colônia não para de morrer. Todo dia morre um, morre dois, morre três. Quando chove, morre mais. Quando faz frio, morre mais. Quando o governo manda verba, morre menos --- mas nunca morre pouco.

O Paciente 1.260 eu conhecia. Não de nome --- ninguém tem nome aqui. Eu conhecia o corpo. Um homem magro, muito magro, com a pele grudada no osso. O olho direito amarelo. Toda vez que eu via ele no pátio, ele estava parado, olhando pro nada. Os enfermeiros chamavam ele de ``o mudo'' porque ele não falava. Sete anos sem falar, diziam. Sete anos calado. E quando ele falou de novo, a primeira palavra foi ``diário''.

Eu pensei: esse homem é um santo ou é um assombrado.

Quando me chamaram pra exumar, eu fui. A cova dele é a 1.260, na vala comum do fundo. A gente cavou. A terra estava fofa, porque a cova era recente. Mas quando a pá bateu no caixão --- não era um caixão, era um lençol, ele foi enterrado num lençol amarrado nas pontas --- mas quando a pá bateu, fez um som diferente. Um som oco, que ecoou. Eu cavei com as mãos.

O senhor quer que eu diga o que eu vi? O senhor não vai acreditar.

O lençol estava intacto. Não tinha barata, não tinha bicho, não tinha nada. O lençol estava branco ainda, como se tivesse sido colocado ontem. Eu desamarrei as pontas. A pele dele estava escura, mas firme. Não tinha cheiro de podre. Tinha um cheiro \textit{adocicado}, como flor de laranjeira, mas mais pesado, que ficou no meu nariz. Eu assoei depois e ainda sentia.

O olho direito --- o amarelo --- ele estava aberto.

O senhor entende? O olho estava \textit{aberto}. O outro, o esquerdo, estava fechado. Mas o amarelo, aberto. Brilhando. Não brilhava como olho vivo, brilhava como vaga-lume --- uma luz fraca, amarela, que vinha de dentro. Eu apaguei a lanterna e vi. Vi a luz no escuro.

E as mãos. As mãos dele estavam cruzadas sobre o peito, segurando um caderno. Um caderno preto, de capa dura, com as bordas desgastadas. As mãos dele estavam em cima do caderno, como quem protege alguma coisa. Quando eu tentei tirar o caderno, os dedos estavam duros. Eu tive que puxar. Fez um som de osso estalando.

O caderno estava quente.

O senhor pode achar que é loucura minha, mas o caderno estava quente. A terra estava fria, o corpo estava frio, mas o caderno estava quente. Como se tivesse acabado de ser escrito.

Eu entreguei o caderno para o doutor. Ele colocou num saco plástico. Eu vi o saco plástico embaçar por dentro, como se o caderno estivesse suando.

Eu não dormi direito desde então.

\vspace{1em}
\hrule
\vspace{1em}

\noindent\textit{Nota do legista:} Durante a exumação, o coveiro apresentou sudorese profusa, taquicardia e tremor nas mãos. Recusou-se a continuar cavando após a abertura do lençol. O restante da exumação foi concluído pelo funcionário administrativo. O caderno (diário) foi acondicionado em saco plástico lacrado e enviado ao IML para análise. O corpo foi recolocado na cova 1.260 e coberto com cal virgem, conforme protocolo do Colônia para valas comuns.

\textit{Observação adicional:} O funcionário administrativo que auxiliou na exumação pediu transferência de setor no dia seguinte. Não informou o motivo. A solicitação foi negada.

\vspace{1em}
\section*{5. DIÁRIO DO PACIENTE 1.260}

\noindent\textit{Transcrição parcial das primeiras páginas. O diário é um caderno de capa preta, 20 cm \texttimes 14 cm, 200 páginas aproximadamente. A capa apresenta manchas de gordura, terra e o que parece ser sangue seco. A lombada está descolada em alguns pontos. As páginas amareladas, com bordas esmaecidas. A caligrafia varia: nas primeiras páginas, é treinada, quase infantil; nas últimas, é errática, violenta, com palavras rasuradas e escritas por cima.}
\vspace{1em}

\subsection*{5.1 Página 1}

\vspace{0.5em}

\noindent\textit{``Este caderno me foi dado pelo Dr. Heitor Oliveira em janeiro de 1979. Ele disse: escreva tudo o que você lembra. Disse que ler me faria bem. Ele não sabe o que está pedindo. O que está escrito aqui não é para ser lido. Mas eu escrevo porque ele pediu. E porque a criatura pede também.''
}

\vspace{0.5em}
\noindent\textit{--- A caligrafia desta primeira página é firme, de alguém que readquiriu o hábito da escrita após décadas de silêncio. Há uma mancha circular na margem inferior, como de copo d'água. O papel está enrugado.}

\vspace{1em}

\subsection*{5.2 Página 2}

\vspace{0.5em}

\noindent\textit{``Meu nome é Joaquim Antônio Correia. Nasci em 1907, em Quixeramobim, Ceará. Fui internado no Hospital Colônia em 1920, aos 13 anos. Motivo da internação: 'alienação mental por epidemia de fome'. Isso é o que está escrito no papel. O que não está escrito é que eu fui internado porque não tinha onde ficar. Porque a seca levou minha família. Porque o governo criou currais para a gente e depois chamou a gente de louco.''
}

\vspace{0.5em}
\noindent\textit{--- Há uma rasura na palavra ``louco''. A tinta atravessou o papel.}

\vspace{1em}

\subsection*{5.3 Página 5}

\vspace{0.5em}

\noindent\textit{``A criatura veio comigo do Ceará. Ela estava na vala dos retirantes, no campo de concentração. Ela estava no olho amarelo do meu pai quando ele morreu sentado. Ela entrou em mim quando eu comi terra da vala. Eu estava com fome. Tanta fome. Comi terra. E a terra continha ela.''
}

\vspace{0.5em}
\noindent\textit{--- A palavra ``fome'' está escrita três vezes, a terceira com tanta pressão que o papel rasgou.}

\vspace{1em}

\subsection*{5.4 Página 8 --- ``A Rasga Mortalha''}

\vspace{0.5em}

\noindent\textit{``A Rasga Mortalha não é uma entidade. É um nome que eu dei para o que vive na vala. As velhas do sertão contam que ela é uma mulher que rasga a mortalha dos mortos para costurar o futuro. Quem ouve o canto dela morre em três dias. Mas eu ouvi e não morri. Eu ouvi e ela entrou.''
}

\vspace{0.5em}

\noindent\textit{``O canto dela é doce. Parece a voz de uma menina cantando acalanto. Mas a letra... a letra fala de ossos. Fala de terra. Fala de um olho amarelo que tudo vê.''
}

\vspace{0.5em}

\noindent\textit{Abaixo, há uma transcrição do canto, em caligrafia trêmula:}

\vspace{0.5em}

\begin{quote}
\itshape
``Rasga a mortalha do menino \\
Rasga a mortalha do pai \\
Rasga a mortalha do mundo \\
O olho amarelo não sai''

\vspace{0.3em}

``A vala é funda, a vala é fria \\
A vala come quem tem fome \\
O olho amarelo é testemunha \\
O olho amarelo não dorme''
\end{quote}

\vspace{0.5em}

\noindent\textit{--- A transcrição termina abruptamente. A palavra ``dorme'' foi escrita e imediatamente coberta por uma mancha de tinta, como se o autor tivesse tentado apagar a palavra depois de escrevê-la. A página 9 está em branco.}

\vspace{1em}

\subsection*{5.5 Páginas 15--22 --- O Mecanismo}

\vspace{0.5em}

\noindent\textit{``A criatura se alimenta de atenção. Quanto mais você pensa nela, mais forte ela fica. Quanto mais você lê sobre ela, mais ela existe. Ela não está no papel. Ela está no espaço entre o papel e o olho. Na hora em que o olho lê a palavra, a criatura salta da página para o cérebro. É por isso que o Dr. Oliveira também começou a ver. Ele leu. Eu mostrei. Eu não devia ter mostrado.''
}

\vspace{0.5em}

\noindent\textit{``Oliveira agora acorda às 3:14 da manhã. Exatamente 3:14. Ele me disse. Ele não sabe que é o horário da vala. É o horário em que o corpo está mais frio. É o horário em que o cérebro está mais aberto. Ela entra quando o corpo está frio.''
}

\vspace{0.5em}

\noindent\textit{``Lúcia não leu. Lúcia só viu o caderno. Ela disse que não ia ler. Mas ela leu. Eu vi no olho dela. O olho direito dela está começando a ficar amarelo.''
}

\vspace{1em}

\subsection*{5.6 Últimas Páginas (ilegíveis)}

\vspace{0.5em}

\noindent\textit{As páginas finais do diário são praticamente ilegíveis. A caligrafia se degrada em rabiscos, espirais, palavras repetidas centenas de vezes. A palavra que mais aparece é ``fome'' --- 47 vezes, segundo contagem preliminar. Depois ``olho'' --- 31 vezes. Depois ``vala'' --- 23 vezes. A última palavra legível, na página 198, é:}

\vspace{0.5em}

\begin{center}
\Large\textbf{\textcolor{olhoamarelo}{\textit{``fome''}}}
\end{center}

\vspace{0.5em}

\noindent\textit{A página 199 está em branco. A página 200 está em branco. Mas na página 200, a luz ultravioleta revelou uma frase escrita a lápis, quase apagada:}

\vspace{0.5em}

\begin{center}
\itshape \small \textcolor{olhoamarelo}{``A criatura está em você agora. Você leu. Você é o 1.263.''}
\end{center}

\vspace{1em}

\section*{6. NOTA DA DR\textsuperscript{A}. LÚCIA MENEZES SOBRE O CICLO DE CONTAMINAÇÃO}

\vspace{0.5em}

\noindent\textit{Anexada ao processo em 20 de junho de 1979. Caligrafia firme, mas com rasuras.}
\vspace{0.5em}

\lettrine{E}{}stou anexando este parecer porque o legista Dr. Álvaro Torres se recusou a incorporá-lo ao laudo oficial. Ele diz que ``não há evidência científica''. Ele tem razão. Não há evidência científica para o que aconteceu. Mas há evidência \textit{testemunhal}. E há o caderno.

O ciclo segue um padrão documentado (embora não explicado):

\begin{enumerate}
    \item \textbf{Contato inicial:} Uma pessoa lê o diário de Joaquim. Não precisa ler inteiro --- uma página, um parágrafo, uma linha é suficiente.
    \item \textbf{Sintomas iniciais (primeiras 24 horas):} A pessoa começa a sentir um cheiro adocicado, sem origem identificável. O cheiro é descrito como ``flor de laranjeira misturada com terra molhada''. A pessoa também começa a acordar às 3:14 da manhã, sem motivo aparente.
    \item \textbf{Fase de incubação (3 a 7 dias):} A pessoa desenvolve um interesse crescente pelo diário. Relê trechos. Sonha com a vala. O sonho é sempre o mesmo: a pessoa está deitada numa vala cheia de corpos, todos com olhos amarelos, todos olhando para ela.
    \item \textbf{Contaminação confirmada (7 a 14 dias):} O olho direito da pessoa começa a apresentar coloração amarelada. Inicialmente sutil, depois mais intensa. A pessoa começa a ouvir o canto da Rasga Mortalha. Ouve como se viesse de dentro da própria cabeça.
    \item \textbf{Fase final (14 a 30 dias):} A pessoa começa a escrever. Escreve sem controle, sem conseguir parar. Escreve as mesmas palavras que Joaquim escreveu. A caligrafia muda, se torna idêntica à de Joaquim. A pessoa se torna o novo elo do ciclo.
\end{enumerate}

\textbf{Paciente 1.260:} Joaquim Antônio Correia (contaminado pela vala dos retirantes em 1915) \\
\textbf{Paciente 1.261:} Dr. Heitor Oliveira (contaminado pela leitura do diário em 1979) \\
\textbf{Paciente 1.262:} Eu, Lúcia Menezes (contaminada pela leitura do diário em 1979) \\
\textbf{Paciente 1.263:} \_\_\_\_\_\_\_\_\_\_\_\_\_\_\_\_\_\_\_\_\_\_\_\_\_\_\_\_\_\_\_\_ (você)

\vspace{0.5em}

\noindent\textit{O nome do Paciente 1.263 está em branco no documento original. Ele deve ser preenchido pelo leitor.}

\vspace{1em}

\noindent\textit{Anotação final de Lúcia, na margem inferior:}

\vspace{0.5em}

\noindent\textit{``Se você está lendo isto, você já é o 1.263. O cheiro já começou? Você já acordou às 3:14? O olho amarelo já apareceu no canto do seu espelho? Não adianta fechar o livro. O ciclo não depende do livro. O ciclo depende de você ter lido. E você leu. Me desculpe. --- Lúcia.''
}

\vspace{2em}

\begin{center}
\fbox{\begin{minipage}{0.92\textwidth}
\centering
\vspace{0.5em}
{\large\textbf{ESCALA DE CONTAMINAÇÃO DO LEITOR}}

\vspace{0.3em}
\small
\textit{Preencha esta tabela ao final da leitura. Some os pontos. Veja o resultado.}

\vspace{0.3em}
\begin{tabularx}{0.88\textwidth}{|c|>{\raggedright\arraybackslash}X|}
\hline
\textbf{Pontos} & \textbf{Diagnóstico} \\ \hline
0--2 & Leitor resiliente. A contaminação não pegou. Ou você está mentindo. \\ \hline
3--5 & Leitor exposto. Sintomas iniciais. Recomenda-se pausa na leitura. \\ \hline
6--8 & Leitor contaminado. O ciclo está em andamento. O olho amarelo aparecerá em breve. \\ \hline
9--10 & \textcolor{olhoamarelo}{Paciente 1.263. Bem-vindo. A criatura agradece sua atenção.} \\ \hline
\end{tabularx}

\vspace{0.3em}
\itshape\textcolor{olhoamarelo}{O ciclo recomeça a cada leitor.}
\vspace{0.5em}
\end{minipage}}
\end{center}

\vspace{2em}

\section*{DESTINAÇÃO DOS MATERIAIS}

\begin{itemize}
    \item \textbf{Corpo:} Liberado para sepultamento em vala comum do Hospital Colônia (Cova 1.260). O corpo foi coberto com cal virgem.
    \item \textbf{Diário:} Retido pelo IML para análise. Anexado a este processo. Permanece sob custódia.
    \item \textbf{Depoimentos:} Gravados em fita cassete e transcritos. Cópias anexadas.
\end{itemize}

\vspace{1em}

\noindent\textit{Nota final do arquivista:}

\vspace{0.5em}

\noindent\textit{O diário do Paciente 1.260 foi roubado do IML de Barbacena em 28 de setembro de 1979. O ladrão nunca foi identificado. Em janeiro de 1980, a Dra. Lúcia Menezes desapareceu. Seu apartamento estava vazio, exceto por uma carta sobre a mesa: ``Voltei para a vala. Não procurem. --- Lúcia''. O Dr. Heitor Oliveira foi encontrado morto em seu consultório em março de 1980. Ao lado do corpo, havia um bilhete: ``A criatura está em mim. Finalmente.''
}

\vspace{0.5em}

\noindent\textit{Algumas páginas do diário foram recuperadas anos depois em diferentes arquivos. As que você leu aqui são as que chegaram até nós. O ciclo continua. O diário está em algum lugar. A criatura está em você.}
\vspace{1em}

\begin{center}
\Large\textit{\textcolor{olhoamarelo}{A página está em branco.}} \\
\Large\textit{\textcolor{olhoamarelo}{A tinta é sua.}}
\end{center}

\vfill

\begin{center}
\fbox{\begin{minipage}{0.7\textwidth}
\centering
\itshape
\small\textcolor{olhoamarelo}{Você leu até aqui.} \\
\small\textcolor{olhoamarelo}{O olho amarelo já está no canto do seu campo de visão.} \\
\small\textcolor{olhoamarelo}{Vire a página.}
\end{minipage}}
\end{center}

\newpage

% Página final --- intencionalmente em branco, exceto por uma linha

\vspace*{\fill}

\begin{center}
\Large\textcolor{olhoamarelo}{\textit{3:14}}
\end{center}

\vspace*{\fill}
